% Development supplement, 11 October 2026.
% The arithmetic dichotomy is unconditional.  The application still assumes
% the proposed local physical moment and the pinned incumbent moment.
\section{Shortened physical scales do not require a new local hypothesis}
\label{sec:physical-scale-shortening}

The local moment window imposes a lower bound on each plain length.
Consequently the sentence ``the varying conductor only shortens the
reflected scale'' does not, by itself, justify applying that moment.
The following dichotomy closes this particular gap. It uses the
\emph{existing} $9/2$ moment on the shorter scales.

Set
\[
 h=\frac{3969737987}{5000000000},\qquad
 \ell=\frac{626701}{3750000},\qquad y=h-2\ell,
 \qquad C=\frac{2249}{500}.
\]
For a row exponent $d\in[h-10^{-3},h+10^{-10}]$ and a free-shift
exponent $t\in[0,10^{-3}]$, normalized relative to $U=Z^d$, the
upper reflected plain length is
\[
 n(d,t)=1-\frac yd+t.
\]
Since $y>0$, this is increasing in both $d$ and $t$. The exact
endpoint values used in the scalar audit are
\[
 n_- =\frac{4998608000}{11894213961},\qquad
 n_+ =\frac{402041377237}{952737117000},
\]
and the nominal selected prime widths lie between
\[
 z_- =\frac{1-2n_+}{C}
      =\frac{74327181263}{2142705776133},\qquad
 z_+ =\frac{1-2n_-}{C}
      =\frac{948498980500}{26750087198289}.
\]
Direct rational comparison gives
\[
 \frac{41}{100}<n_-\le n_+<\frac{211}{500}<\frac{11}{25},
 \qquad
 \frac3{100}<z_-\le z_+<\frac{71}{2000}<\frac1{25}.
\]
All comparisons are strict, so sufficiently small fixed annular,
box, fixed-conductor, and selection-mesh losses fit inside them.
The numerical margins before these losses are approximately
$0.0000143651$ for the upper plain bound $211/500$ and
$0.0000422098$ for the upper slot bound $71/2000$.
In particular those losses must be chosen below these margins;
merely fitting inside the larger local window is insufficient for
the quantitative short-scale reserve below.

\begin{lemma}[A downward scale dichotomy]
Let $n,x_1,x_2,z\ge0$ satisfy
\[
 x_i\le n\le\frac{211}{500},\qquad
 \frac3{100}\le z\le\frac{71}{2000},\qquad
 2n+Cz\le1.
\]
Then either both actual plain lengths $x_i$ belong to
$[41/100,11/25]$ and
\[
 x_1+x_2+Cz\le1,
\]
or
\[
 x_1+x_2+\frac92z\le1-\frac{33}{4000}.
\]
\end{lemma}
\begin{proof}
If both $x_i\ge41/100$, their upper bounds follow from
$211/500<11/25$, and their new capacity follows by comparison
with $2n+Cz$.
If, for example, $x_1\le41/100$, then
\[
 x_1+x_2+\frac92z
 \le\frac{41}{100}+\frac{211}{500}
       +\frac92\frac{71}{2000}
 =1-\frac{33}{4000}.
\]
The other case is symmetric. Equality at the boundary is harmless.
\end{proof}

Here is the precise analytic application. Assume the proposed local
moment has been established for its stated profile and canonical-slot
family, and retain the pinned source's existing $9/2$ moment.
Fix a selected slot family of total width $z$ obeying the lemma,
with the same natural row zeros and eligibility in both estimates.
For any actual pair of reflected scales bounded above by $U^n$, clip
nonempty subunit scales to length zero using the source's bounded
annular dilation rule. Apply the local moment on the rectangle
$x_1,x_2\ge41/100$; on either complementary rectangle apply the
incumbent moment. Both bounds have the same target $U^{1+\epsilon}$.

The original scale-supremum recovery applies separately on these
three rectangles. In logarithmic scale coordinates it bounds the
supremum by integrals of the function and its first derivatives in
each coordinate. These are exactly the four profile pairs appearing
in the two-variable recovery. Their supports remain fixed annuli;
the number of scale boxes is polynomial in $\log U$, and is absorbed
in an arbitrarily small part of $\epsilon$. Thus row-dependent
conductor deficits, downward Euler-restoration scales, and the low
annuli of the reflected Schwartz profile do not force a local-moment
claim below $41/100$. The upper annular tail is first removed with
the original rapid-decay estimate.

The algebra relating this to the actual numerator is the pinned
source's B0.42--B0.44: write a monomial shift as $h_0h_r$, with
$h_r$ charged to the repeated-prime conductor deficit. Its reflected
scale is bounded by a fixed constant times $(U/Y)h_0$; natural Euler
restoration only shortens that scale. For a fixed $h_0$ dyad this
is the upper bound $n(d,t)$ above. The same combined conductor must
be used in this bound and in the coefficient damping. No factor of
the conductor deficit may be charged twice. The argument introduces
no further coefficient mass and therefore preserves the source's
weighted free-shift mass $\sum |c_u(h_0,h_r)|h_0^{\delta/2}$.

\paragraph{Short actual prime supply.}
There is a second possible shortening: the actual selected prime width
$z$ can be smaller than the nominal selection capacity
$z_A=(1-2n)/C$. Error slots carry zero gain, and even a main slot
may carry zero gain. No lower bound on the positive-gain supply is
required. If $z\le3/100$, the old moment applies directly, since
\[
 x_1+x_2+\frac92z
 \le 2\frac{211}{500}+\frac92\frac3{100}
 =1-\frac{21}{1000}.
\]
Thus every actual selected width $0\le z\le z_A\le71/2000$ is
covered: use this old-moment branch for $z\le3/100$, and the
preceding downward-scale dichotomy otherwise. The common worst
reserve in either old-moment branch is $33/4000$.

The selection gain is still the gain at the \emph{target capacity}
$z_A$, not merely at the smaller actual width $z$. To see this
from the literal source, write $w_i$ for every original slot width,
$L=\sum_iw_i$ for their total, and $g_i\in[0,\delta/2]$ for the
main-masked amplitude-bin gains. The error gains are zero but their
widths are retained in $L$. Set
\[
 q=\frac{\sum_iw_i g_i}{L}.
\]
The source theorem \path{zero_error_whole_slots} assumes only
$0\le z_A\le L$, $w_i\le\mathrm{mesh}$, and the above gain bounds.
It returns a subset $I$ of the positive-gain main slots satisfying
\[
 z=\sum_{i\in I}w_i\le z_A,\qquad
 \sum_{i\in I}w_i g_i\ge qz_A-\frac\delta2\mathrm{mesh},
 \qquad
 \left|\prod_{i\in I}Q_i\right|^2
 \ge U^{2qz_A-\delta\mathrm{mesh}}.
\]
There is no premise that the positive-gain main supply is at least
$z_A$. If it is exhausted first, the selected weighted gain is
$qL\ge qz_A$, which is stronger than the stated rounded bound.
Zero-gain main factors must not be reinserted into this lower-bounded
product; such factors could be arbitrarily small or zero.

The source's \path{source_class_main_selection} constructs the
amplitude classes on all original slots and uses exactly this masked
mean. Consequently the same $I$, $q$, and lower bound work uniformly
on its class of rows. Its full-main \emph{upper} bound is unchanged.
To substitute the local capacity for its old capacity one invokes
\path{zero_error_whole_slots} at $z_A$ directly, after checking
the total-original-supply inequality $z_A\le L=\ell/d$. This is the
scalar audit's capacity gate; it does not assert a fictitious
positive-main supply lower bound.

Combining this selection with either applicable physical moment
gives the same denominator gain $U^{2qz_A}$, up to the fixed mesh
loss, whether the actual width is in the local window or below it.
The old-moment branch therefore preserves the desired weighted
numerator saving. This argument supplies neither the new local
moment nor its literal initial-state adapter; it closes the
downward-scale and short-selected-supply parts conditional on those
inputs and the pinned incumbent moment.

The six elementary statements in \path{ScaleShortening.lean}
check the rational upper edges, both strict incumbent reserves,
the two dichotomies, and the exhausted-supply affine gain. Their
scope is arithmetic; the actual selection and spike identity are
the source theorems identified above.

\paragraph{All larger free shifts, including zero capacity.}
For completeness the restriction $t\le10^{-3}$ does not discard the
remaining coefficient dyads. Put $b=2y/d-1>0$, $t_0=10^{-3}$,
and
\[
 z_{\mathrm{new}}(0)=\frac bC,\qquad
 z_{\mathrm{old}}(t)=\max\left\{0,\frac{b-2t}{9/2}\right\}.
\]
The actual incumbent capacity uses $6\kappa\le9/2$, so this is a
lower admissible capacity. The weighted coefficient mass in B0.44
pays $U^{-\delta t/2}$. Thus it suffices to prove
\[
 \frac{\delta t}{2}+\frac{qz_{\mathrm{old}}(t)}2
 \ge\frac{qz_{\mathrm{new}}(0)}2.
 \tag{SS}
\]
Write
\[
 S_b(t)=\frac t2-\frac14\bigl(b/C-z_{\mathrm{old}}(t)\bigr).
\]
This function is continuous at $t=b/2$. Its derivative is $7/18$
where the old capacity is positive and $1/2$ where it is zero.
For the audited near-endpoint rows,
$b\le b_{\max}=1-2n_-$ and $b>2t_0$.
The exact rational audit gives
\[
 S_b(t_0)\ge
 \frac{t_0}{2}-\frac14
   \left(\frac{b_{\max}}C-
         \frac{b_{\max}-2t_0}{9/2}\right)
 >\frac3{10000}.
\]
It follows that $S_b(t)>0$ for every $t\ge t_0$, with no exception
at the clipping point. Since $0\le q\le\delta/2$ and
$z_{\mathrm{new}}(0)-z_{\mathrm{old}}(t)\ge0$, the left side of
(SS) minus its right side is at least $\delta S_b(t)\ge0$.
This proves the full large-shift payment, including the empty selected
family at zero capacity. That family uses the incumbent unrestricted
zero-slot moment B3.40, not a moment with a hidden requirement
$x_1+x_2\le1$. In the source definition \path{ZeroAt}, the
positive scales are required only to satisfy $X_i\le Z^L$; its
conclusion is $Z^{s.\mathrm{width}+\epsilon}$. The actual physical
specialization \path{finite_pair_zeroAt} in
\path{Cycle25/Energy/Physical/CertifiedZero.lean} gives
$U^{1+\delta+\epsilon}$ for $X_i\le U^2$.
Its \path{certified_reflected_zero_fourier_moment}
retains both independent Fourier heights.
Finally \path{certified_original_numerator_fourth} in
\path{Cycle25/Energy/Physical/OriginalFourth.lean} accepts
the exact disjunction
\[
 2(m+\xi)+6\kappa\sum_{i\in I}w_i\le1
 \quad\text{or}\quad I=\varnothing,
 \qquad m+\xi\le2,
\]
and gives the same $U^{1+o(1)}$ fourth-moment bound in either case.
There is consequently no extra positive-part length cost when the
capacity vanishes. The larger $\max(1,2m)$ bound in the unmarked
density adapter allows its finitely many exceptional rows; it is
not the eligible physical zero-slot estimate being used here.

For $0\le t\le t_0$ the local capacity is
$z_{\mathrm{new}}(t)=(b-2t)/C>0$.
Here the same conclusion follows directly from
\[
 \frac{\delta t}{2}+
 \frac q2\bigl(z_{\mathrm{new}}(t)-z_{\mathrm{new}}(0)\bigr)
 =t\left(\frac\delta2-\frac qC\right)\ge0.
\]
The original support cutoff restricts $t\le m+o(1)$, as in B0.44;
all scales therefore stay in bounded geometric ranges. The new
argument requires no new uniform estimate for unbounded plain
lengths. Row exponents below the near-endpoint interval use the
separate incumbent row-endpoint certificates, as specified in the
global endpoint assembly.
